\begin{exercisechunk}
\item \textbf{Attention as a modern Hopfield update.}\suppmark
\label[exercise]{ex:l14-hopfield-attention}
\exerciseprofile{2}{2}{2}{\exproof\quad\excalculation\quad\exinterpretation}{%
Derive the energy update underlying softmax attention and quantify its approach
to hard associative lookup.}
Earlier exercises analyze how score ranges control softmax concentration and
retrieval error. The new result here is the energy interpretation: when keys
are also values, one softmax-attention update is one gradient step for a
modern Hopfield energy.

Let \(M\ge2\), let \(k_1,\ldots,k_M\in\R^d\), let \(\beta>0\), and define
\[
 \alpha_i(q)=\frac{e^{\beta k_i^\top q}}
 {\sum_{j=1}^M e^{\beta k_j^\top q}},
 \qquad
 E_\beta(q)=\frac12\|q\|_2^2-
 \frac1\beta\log\left(\sum_{i=1}^M e^{\beta k_i^\top q}\right).
\]
\begin{enumerate}[(a),beginpenalty=10000]
\item Differentiate \(E_\beta\) and prove
\[
 \nabla E_\beta(q)=q-\sum_{i=1}^M\alpha_i(q)k_i.
\]
\item Show that one gradient step of length one gives
\[
 q^+=q-\nabla E_\beta(q)
 =\sum_{i=1}^M\alpha_i(q)k_i.
\]
Identify the queries, keys, and values in the corresponding one-query
attention operation. Explain why a general attention output
\(\sum_i\alpha_i(q)v_i\) with arbitrary \(v_i\) need not be a gradient step
in the query space \(\R^d\).
\item Suppose \(i^\star\) uniquely maximizes \(k_i^\top q\), define
\[
 \Delta=k_{i^\star}^\top q-\max_{i\ne i^\star}k_i^\top q>0,
\]
and assume \(\|k_i\|_2\le B\) for every \(i\). You may use
\[
 1-\alpha_{i^\star}(q)
 \le (M-1)e^{-\beta\Delta}.
\]
Prove
\[
 \|q^+-k_{i^\star}\|_2
 \le2B(M-1)e^{-\beta\Delta}.
\]
For \(0<\varepsilon<2B(M-1)\), give a sufficient lower bound on \(\beta\)
for this error to be at most \(\varepsilon\).
\item Explain which limit produces hard associative lookup and identify the
roles of the score margin \(\Delta\), the number of memories \(M\), and the
norm bound \(B\) in the quantitative approximation.
\end{enumerate}
\end{exercisechunk}
